\labsessionheader{2026-09-30}{Batch scrub-guard compiler queries}%
  {15663460 + session changes}{n/a}%
  {docs/audits-2026-09-30/scrub\_query\_batching.json}

\subsection*{Changes}
\bd{test\_C\_libtooling-b1l3}: the root probe now asks previously encountered
non-template record identities and member-type questions in one compiler
translation unit. Only answers from that pass's AST and compiler flags are
reused. Guesses contain names and member accesses, never previous AST nodes or
verdicts. Template instances are excluded from speculation. Missing answers
use the existing probes; errors outside speculative lines also fall back.

The new 96-line \texttt{tools/scrub\_guard\_batch.py} shares the existing
alias-to-compiler-ID resolution with ordinary probes. Alias and member-query
answers come only from the helper namespaces; definitions are selected by AST
ID from the full dump. The file loader also supports reviewer harnesses that
import the guard by path without registering it in \texttt{sys.modules}.
No mutation corpus entries were removed.

\subsection*{Focused measurements}
Ran \texttt{python3 /tmp/cq\_guard\_profile.py} before and after batching,
with \texttt{set\_jobs(2)} and a raw-pointer-field check followed by a pristine
check. The pristine API check runs four configuration/PCH passes: it took
19.62 seconds and 19 compiler invocations before, versus 15.14 seconds and
15 invocations after, with the same clean result and 78 members. The cold
raw-field check took 7.45 seconds before and 7.33 seconds after; both reported
one raw-pointer problem and 79 members. These are single samples.

\subsection*{Verification}
\texttt{python3 /tmp/cq\_batch\_controls.py} checked fallback on an empty AST
dump, selection of a record by compiler ID when printed names collide,
selection of a template owner's identity instead of its argument's identity,
and preservation of the uninstantiated-template verdict. A real no-PCH guard
run through an unregistered importlib module was clean over 78 members.
The full \texttt{ctest --test-dir build -j6 --output-on-failure} run passed
585/585 entries in 375.65 seconds (375.86 seconds external wall time), versus
492.98 seconds (493.02 external) in session 032: approximately 24 percent less
wall time in these two runs. Both had existing stage caches enabled; these are
single consecutive measurements, not a repeated controlled benchmark.
The three mutation shards took 375.60, 339.47 and 348.59 seconds versus
492.93, 437.84 and 446.40 previously. Every shard reported zero failures,
covering the same 157 mutants and 19 expected-clean corpus entries.

The required \texttt{cmake --build build --target cq-opt-equivalence-full}
sweep compared 2,489 fixtures with zero mismatches in 118.28 seconds external
wall time, versus 147.75 seconds previously, despite unchanged compiler code.
Runtime-condition variation therefore contributes to the observed full-suite
reduction; the compiler-invocation count is the direct evidence of batching's
work reduction. \texttt{python3 tools/update\_code\_map.py --check} and
\texttt{git diff --check} passed.

\subsection*{Follow-up}
\bd{test\_C\_libtooling-2n7r} tracks the remaining decomposition of the
existing large guard tool. This change extracts only compiler identity and
batching concerns.
